\rok{Новембар 2024.}{G}

Први практични тест из Алгоритама и структура података

\zadatak{1}{Бинарно претраживање}
Написати \textit{search} методу класе \texttt{BinarySearch}.

\textbf{Додатне напомене за решавање задатка:}
\begin{itemize}
    \item Улаз је сортиран низ целих бројева.
    \item Излаз је \texttt{true} или \texttt{false}, у зависности да ли је број пронађен.
    \item У шаблону се налази класа \texttt{SortedArrayGenerator} коју треба искористити за генерисање сортираног низа случајних бројева.
    \item Креирати класу \texttt{BinarySearch} и у оквиру ње генерисати методу:
    \begin{Verbatim}[fontsize=\footnotesize, numbers=left, xleftmargin=10mm]
public static bool search(long[] arr, int size, long value)
    \end{Verbatim}
    \item Тестирати у оквиру \texttt{Main} методе.
\end{itemize}



\zadatak{2}{Промена редоследа n елемената са краја реда, користећи стек}
Имплементирати функцију која за задати ред целих бројева обрће редослед само последњих n елемената реда, користећи помоћни стек. Преостали елементи реда остају у истом редоследу.

\textbf{Спецификација функције:}
\begin{itemize}
    \item Функција прима број n, где је n број елемената са краја реда које треба обрнути.
    \item Функција треба користити стек (\textit{stack}) као помоћну структуру.
    \item Елементи пре последњих n треба да остану у оригиналном редоследу.
\end{itemize}

\textbf{Додатни захтеви за решавање задатка:}
\begin{itemize}
    \item Алгоритам писати у оквиру класе \texttt{Queue}.
    \item Захтеван потпис методе:
    \begin{Verbatim}[fontsize=\footnotesize, numbers=left, xleftmargin=10mm]
public void ReverseLastNQueueElements (int n)
    \end{Verbatim}
\end{itemize}

\textbf{Примери:}
\begin{itemize}
    \item \textbf{Пример 1:}
    \begin{itemize}
        \item Оригинални ред: 1 2 3 4 5 6 7, $n = 4$
        \item Ред након обртања последњих 4 елемената: 1 2 3 7 6 5 4
    \end{itemize}
    \item \textbf{Пример 2:}
    \begin{itemize}
        \item Оригинални ред: 5 3 2 1 8, $n = 5$
        \item Ред након обртања последњих 5 елемената: 8 1 2 3 5
    \end{itemize}
    \item \textbf{Пример 3:}
    \begin{itemize}
        \item Оригинални ред: 1 2 3 4 5 6, $n = 7$
        \item Неважећа вредност за n.
    \end{itemize}
\end{itemize}



\zadatak{3}{Најдужи сегмент са непарним и парним чворовима}
Чворови у повезаној листи могу формирати сегмент са наизменично непарним и парним вредностима ако задовољавају следеће услове:
\begin{enumerate}
    \item Сваки чвор у сегменту има вредност која је различита по парности од вредности претходног чвора.
    \item Сегмент мора имати најмање 3 чвора.
\end{enumerate}

Задатак је пронаћи и вратити:
\begin{enumerate}
    \item Дужину најдужег сегмента са наизменично непарним и парним чворовима.
    \item Почетну вредност сегмента.
    \item Крајњу вредност сегмента.
\end{enumerate}

Ако ниједан такав сегмент не постоји, треба вратити низ [-1, -1, -1].

\textbf{Додатни захтеви:}
\begin{itemize}
    \item Решавај овај задатак помоћу структуре листе (\textit{LinkedList}) из шаблона задатка.
    \item Имплементирај га у оквиру методе:
    \begin{Verbatim}[fontsize=\footnotesize, numbers=left, xleftmargin=10mm]
public int[] findLongestSegment()
    \end{Verbatim}
\end{itemize}

\textbf{Примери:}
\begin{itemize}
    \item \textbf{Пример 1:}
    \begin{itemize}
        \item \textbf{Улаз:} листа: $1 \rightarrow 3 \rightarrow 5 \rightarrow 7$
        \item \textbf{Излаз:} [-1, -1, -1]
        \item Објашњење: Нема наизменичних вредности по парности, па не постоји валидан сегмент.
    \end{itemize}
    \item \textbf{Пример 2:}
    \begin{itemize}
        \item \textbf{Улаз:} листа: $2 \rightarrow 3 \rightarrow 4 \rightarrow 5 \rightarrow 6 \rightarrow 7$
        \item \textbf{Излаз:} [6, 2, 7]
        \item Објашњење: Најдужи сегмент је $2 \rightarrow 3 \rightarrow 4 \rightarrow 5 \rightarrow 6 \rightarrow 7$, са дужином 6. Почиње на вредности 2 и завршава се на вредности 7.
    \end{itemize}
\end{itemize}

\sablon{Шаблон првог теста новембар 2024. група G}{2024/Novembar/GrupaG/AISP2024_Test1_GrupaG.linq}
